\setcounter{equation}{0}
\setcounter{Q}{0}

%TGM1, pg51
Le diagramme suivant représente le diagramme d'équilibre binaire fer-carbone.

\begin{center}
\begin{minipage}{0.53\textwidth}
\centering
\includegraphics[width=\textwidth]{/home/jeff/Documents/IUTOrsay/1-Cours/4p13-Alliages/Cours-TD/n-3-Lecture&Utilisation_Binaires&Ternaires/Chp3-Ex/Aciers-Ex02/figures/Fig01}
\end{minipage}\hfill\begin{minipage}{0.43\textwidth}
\captionof{figure}{Un diagramme binaire fer-carbone}
\end{minipage}
\end{center}

\Q Quel est le type de ce diagramme (stable ou métastable)
\Q Indexer le diagramme en précisant le nombre de phase de chaque domaine.
\Q Donner les coordonnées des points particuliers dans ce diagramme, en précisant pour
chacun le type de la transformation et l'équation d'équilibre.
\Q Soit l'alliage à \SI{1.5}{\percent} de carbone.
\begin{enumerate}
\item Calculer la proportion en ferrite et cémentite dans cet alliage.

\item En considérant que la perlite forment des grains de composition fixe de \SI{0.77}{\percent} en carbone. Calculer le pourcentage de grain de perlite et de cémentite dans le mélange.

\item Donner le nom de cet alliage.
\end{enumerate}